\documentclass[a4paper,11pt]{ltjsarticle}
\usepackage{amsmath,amssymb,amsthm,bm,booktabs,array,longtable,mathtools}
\usepackage{hyperref}
\hypersetup{unicode=true,hidelinks}
\setlength{\emergencystretch}{2em}
\newtheorem{theorem}{定理}
\newtheorem{proposition}[theorem]{命題}
\newtheorem{lemma}[theorem]{補題}
\newtheorem{corollary}[theorem]{系}
\theoremstyle{definition}
\newtheorem{definition}[theorem]{定義}
\newtheorem{remark}[theorem]{注意}
\newcommand{\tr}{\operatorname{tr}}
\newcommand{\Ric}{\operatorname{Ric}}
\newcommand{\eps}{\varepsilon}

\title{ハード・ジャイロスコピック重力波メモリーに対応する固定窓付き最小ヌル曲率作用：鋭いBel--Robinson下界}
\author{}
\date{}

\begin{document}
\maketitle

\begin{abstract}
本稿では、弱い重力波が生むジャイロスコピック重力メモリーのハード／ヘリシティ成分と、同じ局所放射域パッチの放射データを横切る固定した補助的ヌル Jacobi プローブ上の Bel--Robinson 曲率作用との間に鋭い下界を導く。対象は四次元、放射域の先頭次数、標準的な初期シアーなし Bondi フレーム、入出力定常条件を満たす線形化 TT/Bondi 放射データのクラスに限定する。ヌル Jacobi 散乱量 $\omega_\gamma$ と時間的世界線上のジャイロ向き角 $\Delta\Psi_{\rm H}$ は別個の観測量であるが、同じ局所放射データ上の二次汎関数として
\[
\omega_\gamma^{(2)}=-2\Delta\Psi_{\rm H}^{(2)},\qquad
|\omega_\gamma^{(2)}|=2|\Delta\Psi_{\rm H}^{(2)}|
\]
が成り立つ。

入出力の定常性は二偏光の潮汐制御に平均ゼロ条件を課す。このため物理的最小作用問題を
\[
\mathcal U_0=\left\{f\in L^2(0,1):\int_0^1 f(t)\,dt=0\right\},\qquad
K_{\rm stat}=P_0K_0P_0,
\]
\[
(K_0f)(t)=\frac12\int_0^1(t-s)|t-s|f(s)\,ds
\]
上の圧縮スペクトル問題として解く。ここで $P_0$ は $\mathcal U_0$ への直交射影である。定常 Jacobi 最小作用関数は
\[
\kappa_{\rm stat}(\omega)=\frac{|\omega|}{\|K_{\rm stat}\|}+O(\omega^2)
\]
を満たし、その係数は厳密終端回復によって鋭い。弱振幅族に対して
\[
Q_\gamma(\eps)=\eps^2Q_\gamma^{(2)}+O(\eps^3),\qquad
\Delta\Psi_{\rm H}(\eps)=\eps^2\Delta\Psi_{\rm H}^{(2)}+O(\eps^3)
\]
と二次摂動\emph{係数}を定義する。この leading-order 放射データ問題は純粋な二次変分問題であり、
\[
\boxed{\kappa_{\rm H}^{\rm lead}(\psi)=\frac{2}{\|K_{\rm stat}\|}|\psi|}
\]
と厳密に解ける。ここで \(\psi=\Delta\Psi_{\rm H}^{(2)}\) は full finite-amplitude memory ではなく、その二次摂動係数であり、\(\kappa_{\rm H}^{\rm lead}\) は \(Q_\gamma^{(2)}\) を最小化する。Gauss--Legendre Nystr\"om 法による検算では
\[
\|K_{\rm stat}\|=0.0281140680170\ldots,\qquad
\frac{2}{\|K_{\rm stat}\|}=71.13876223\ldots
\]
を得る。十進値は数値検証であり、定理の係数は作用素ノルムによって厳密に定義する。

ここで
\[
Q_\gamma=L^3\int_\gamma B(k,k,k,k)\,d\lambda
\]
は、固定した放射横断ヌル区間 $\gamma$ に付随する無次元の windowed Bel--Robinson 曲率作用である。これはジャイロの時間的世界線に沿う曝露量ではなく、また Bondi エネルギー、Isaacson エネルギー、放射エネルギー、検出器エネルギーでもない。固定区間のアフィン再パラメータ化には不変だが、区間端点そのものを変えれば別の最小化問題になる。また、異なる放射横断ヌル方向や異なるスクリーン同定を選べば別の補助プローブ問題であり、本稿は $Q_\gamma$ を方向独立な固有曝露量とはみなさない。本稿の主定理が最小化するのは、この固定ヌル窓上の量の弱振幅展開における先頭非零係数 $Q_\gamma^{(2)}$ である。本稿の主張はハード／ヘリシティ成分、弱振幅、放射域の先頭次数、初期シアーなしフレームに限定し、最終シアーは零に固定せず一定の displacement-memory 成分を許す。ジャイロスコピック・メモリー全体、スピン・メモリー、全 BMS フレームで不変な主張、有限振幅一様定理は主張しない。また、ここで構成する任意の回復波形が完全な漸近平坦真空 Einstein 時空へ大域的に埋め込めることも証明しない。
\end{abstract}

\section{はじめに}
重力波メモリーとは、一過性の重力放射が通過した後にも残る持続的な効果の総称である。変位メモリーに加え、スピン・メモリー、重心メモリー、ホロノミー、曲線偏差や回転する試験粒子に基づく持続的観測量など、複数の観測量が研究されている\cite{Flanagan2019,Flanagan2020,GrantNichols2022,SerajNeogi2023,Harte2025}。本稿が対象とするジャイロスコピック重力メモリーは、遠方の自由落下ジャイロが重力波通過後に獲得する向きの変化であり、先頭の $r^{-2}$ 項はソフト成分とハード成分に分かれる。既存研究ではソフト成分はスピン・メモリーと一致し、ハード成分は Bondi シアーに二次で、局所的な重力波ヘリシティおよび放射位相空間上の局所的な電気・磁気双対性の生成子として解釈される\cite{SerajOblak2023,FayeSeraj2025}。したがって、本稿はジャイロスコピック・メモリーやヘリシティを新たに発見するものではない。

本稿の問いは逆向きである。所定のハード・ジャイロスコピック・メモリーを持つ放射データに対し、その同じ波形を横切る固定ヌル窓上の Bel--Robinson 曲率作用はどこまで小さくできるかを問う。ここで資源として用いる $Q_\gamma$ は、ジャイロ世界線上の作用ではなく、補助的な放射横断ヌル測地線区間 $\gamma$ に沿う四重ヌル縮約の積分である。exact Jacobi 問題では $Q_\gamma$ 自体を用い、Bondi/TT 放射データとの比較では弱振幅展開の二次係数 $Q_\gamma^{(2)}$ を用いる。本稿で「鋭い」とは、単に一つの下界定数を与えるだけでなく、その定数より大きい普遍係数では不等式を維持できず、許容回復族によってその下界を任意の精度で回復できることを意味する。

本研究の数学的核は、二偏光のトレースレスなヌル Jacobi プロファイルに対するコンパクト反自己共役作用素
\[
(K_0f)(t)=\frac12\int_0^1(t-s)|t-s|f(s)\,ds
\]
である。無拘束の非線形 Jacobi メモリー・チャネルに対しては $\|K_0\|^{-1}$ が厳密な小目標の最良係数となる。しかし、物理的な入出力定常条件は潮汐制御に平均ゼロ条件を課すため、ハード・ジャイロスコピック・メモリー問題の最良係数は全 $L^2$ 空間のノルムではなく、平均ゼロ部分空間への圧縮
\[
K_{\rm stat}=P_0K_0P_0
\]
のノルムで決まる。本稿ではこの制約付き最小作用問題を解析的に解き、標準的な初期シアーなし Bondi フレームにおける Jacobi チャネルとハード・ジャイロスコピック・メモリーを結ぶ係数 2 の対応関係と組み合わせる。その結果、先頭次数のハード・メモリー二次係数を固定する限定放射データ問題に対して
\[
\kappa_{\rm H}^{\rm lead}(\psi)=\frac{2}{\|K_{\rm stat}\|}|\psi|
\]
という、$Q_\gamma^{(2)}$ に対する厳密で鋭い最小作用則を得る。exact nonlinear Jacobi endpoint の値関数は別に \(O(\omega^2)\) の補正を伴う小目標漸近則であり、両者を区別する。

既存研究との差は次のように限定する。Jacobi 伝播作用素、平面波メモリーテンソル、Bel--Robinson テンソル、ジャイロスコピック・メモリー、$N_{AB}\widetilde C^{AB}$ 型ハード汎関数、ヘリシティ／双対性の解釈は既知である\cite{Flanagan2020,Harte2025,SerajOblak2023,FayeSeraj2025}。本稿で新たに導く結果は、これらの既知構造に対して、同じ放射データ上のハード・ジャイロスコピック／ヘリシティ汎関数と固定ヌル窓上の Bel--Robinson 曲率作用を制約付き逆最適化として結び、その最良係数を確定する点にある。最終定理の同時条件に対する文献検索では対応する先行定理を特定できなかったが、この検索結果を歴史的優先権の証明とはみなさない。

入出力定常条件を課すと、無拘束問題の係数 $2/\|K_0\|=21.96874242\ldots$ は最良ではなくなる。定常条件から $\alpha,\beta\in\mathcal U_0$ が従い、この閉部分空間上で厳密終端展開とコンパクト作用素による回復構成をやり直すと、最良係数は $2/\|K_{\rm stat}\|$ となる。したがって旧係数は有効な弱い下界としては残るが、最良係数ではない。本稿の「鋭い」は、この制限作用素定理によって担保される。

\section{幾何学的・物理的設定}
\subsection{Bondi フレームとハード・ジャイロスコピック・メモリー}
Bondi 座標を $(u,r,\theta^A)$ とし、
\[
g_{AB}=r^2\gamma_{AB}+rC_{AB}+O(1),\qquad N_{AB}=\partial_u C_{AB}
\]
とする。本稿ではハード／ソフト分解の基準として、標準的な初期シアーなし真空フレーム
\begin{equation}
C_{AB}(u=-\infty)=0
\label{eq:early-shear-free}
\end{equation}
を固定する。最終シアー $\Delta C_{AB}$ は非零でよく、通常の変位メモリーを排除しない。

Seraj--Oblak および Faye--Seraj の規約に合わせ、$\widetilde C_{AB}=\epsilon_{CA}C^C{}_B$ とする。ジャイロスコープ歳差のハード成分は
\begin{equation}
\widehat\Omega_{\rm H}=-\frac18N_{AB}\widetilde C^{AB}
\label{eq:hard-precession}
\end{equation}
であり、物理的な向き角は放射域で
\begin{equation}
\Delta\Psi_{\rm H}=\frac{1}{r^2}\bar\Phi_{\rm H}+O(r^{-3}),\qquad
\bar\Phi_{\rm H}=-\frac18\int du\,N_{AB}\widetilde C^{AB}.
\label{eq:hard-memory}
\end{equation}
本稿では $\Delta\Psi_{\rm H}$ をハード／ヘリシティ・ジャイロスコピック・メモリーと呼ぶ。これはジャイロスコピック・メモリー全体でもスピン・メモリーでもない\cite{SerajOblak2023,FayeSeraj2025}。

\subsection{ヌル Jacobi 系}
四次元 Lorentz 時空で、放射を横切る補助的なヌル測地線区間 $\gamma:[\lambda_0,\lambda_1]\to M$ を考え、$\lambda$ をアフィン・パラメータ、$k=d\gamma/d\lambda$, $L=\lambda_1-\lambda_0>0$ とする。$\gamma$ はジャイロスコープの時間的世界線ではない。また、放射域で outgoing wave と共伝播する $u=\mathrm{const.}$ のヌル生成線も対象外とし、先頭放射曲率を横切る $du/d\lambda\neq0$ のヌル線を選ぶ。本稿ではさらに、$r\to\infty$ で $L$ を固定して $L/r\to0$ とする局所放射域パッチを用い、$\gamma$ の角方向・半径方向のドリフトによる効果を先頭 $1/r$ 次数の外へ送る。スクリーン基底は、この局所パッチで Bondi/TT 偏光二脚の横方向成分と一致するよう適合させる。これらは主定理の幾何学的仮定であり、任意の有限 $r$ のヌル方向に対する方向独立な主張ではない。

先頭の漸近平坦・局所平面波近似では $du/d\lambda$ は一定であり、正のアフィン自由度を用いて $du/d\lambda=1$ と規格化し、選択した stationary-to-stationary 窓の retarded-time 幅と $L$ を一致させる。$e_A$ を平行移動された正規直交スクリーン基底とし、
\begin{equation}
T_{AB}(\lambda)=-R(e_A,k,e_B,k)
\label{eq:tidal}
\end{equation}
とする。$\Ric(k,k)=0$ なら $T$ は対称トレースレスな $2\times2$ 行列である。

$t=(\lambda-\lambda_0)/L$ と無次元化し、
\[
A(t)=L^2T(\lambda_0+Lt)
\]
を定める。Jacobi 基本行列 $\Phi$ は
\begin{equation}
\frac{d\Phi}{dt}=
\begin{pmatrix}0&I_2\\ A(t)&0\end{pmatrix}\Phi,
\qquad \Phi(0)=I_4
\label{eq:jacobi}
\end{equation}
を満たす。自由伝播行列を $F_0(s)=\begin{psmallmatrix}I_2&sI_2\\0&I_2\end{psmallmatrix}$ として
\begin{equation}
S_\gamma=F_0(-1)\Phi(1)
\label{eq:S}
\end{equation}
を、平坦時空での自由伝播を除いた散乱行列とする。ブロック分解 $S_\gamma=(S_{ij})$ に対して
\begin{equation}
\operatorname{skew}(S_{11}+S_{22})=\omega_\gamma J_2,
\qquad J_2=\begin{pmatrix}0&-1\\1&0\end{pmatrix}
\label{eq:omega-def}
\end{equation}
で定まる実スカラー $\omega_\gamma$ を非線形 Jacobi メモリー／散乱チャネルと呼ぶ。これは平面波メモリーテンソル表示では、対角 Jacobi ブロックの反対称結合に対応する\cite{Flanagan2020,Harte2025}。

\subsection{Bel--Robinson 曲率曝露}
計量の符号を $(-+++)$ とし、Bel--Robinson テンソルを
\begin{equation}
B_{abcd}=\frac14\left(C_{aecf}C_b{}^e{}_d{}^f+{}^\star C_{aecf}\,{}^\star C_b{}^e{}_d{}^f\right)
\label{eq:BR-def}
\end{equation}
と規格化する。真空のヌル・スクリーンでは
\begin{equation}
B(k,k,k,k)=|\Psi_0|^2=\frac12\tr(T^2).
\label{eq:BR-tidal}
\end{equation}
そこで
\begin{equation}
Q_\gamma:=L^3\int_{\lambda_0}^{\lambda_1}B(k,k,k,k)\,d\lambda
=\frac12\int_0^1\tr(A(t)^2)\,dt
\label{eq:Q}
\end{equation}
と定義する。$\lambda'=a\lambda+b$（$a>0$）という正のアフィン再パラメータ化では
$L'=aL$, $k'=a^{-1}k$, $d\lambda'=a\,d\lambda$,
$B(k',k',k',k')=a^{-4}B(k,k,k,k)$ なので、
$L'^3d\lambda'\,B(k'^4)=L^3d\lambda\,B(k^4)$ となる。従って $du/d\lambda=1$ は固定した放射横断ヌル区間のアフィン同値類から代表を選ぶ規格化であり、$Q_\gamma$ 自体のアフィン不変性を壊さない。ただし、区間端点を動かして窓そのものを変える、または別の放射横断ヌル方向・別のスクリーン同定を選ぶ場合は別の補助プローブ問題であり、$Q_\gamma$ は window-independent / direction-independent な量ではない。

\begin{remark}
$Q_\gamma$ は Weyl 曲率の二次量に由来する無次元の窓付きヌル作用であり、ジャイロ世界線上の曝露量、Bondi 質量減少、Isaacson 有効応力エネルギーテンソル、放射エネルギー、検出器エネルギーのいずれとも同一視しない。
\end{remark}

\section{Jacobi メモリーと最小作用問題}
STF 基底
\[
\sigma_3=\begin{pmatrix}1&0\\0&-1\end{pmatrix},\qquad
\sigma_1=\begin{pmatrix}0&1\\1&0\end{pmatrix}
\]
を用いて
\begin{equation}
A(t)=\alpha(t)\sigma_3+\beta(t)\sigma_1,
\qquad Q_\gamma=\|\alpha\|_2^2+\|\beta\|_2^2.
\label{eq:control}
\end{equation}

\begin{definition}
\begin{equation}
(K_0f)(t)=\frac12\int_0^1(t-s)|t-s|f(s)\,ds,
\qquad s_0:=\|K_0\|,
\label{eq:K0}
\end{equation}
とする。核は反対称なので $K_0^*=-K_0$、かつ Hilbert--Schmidt 作用素なのでコンパクトである。
\end{definition}

\begin{lemma}[厳密終端写像の小振幅展開]
\label{lem:endpoint-expansion}
十分小さい $A$ に対して
\begin{equation}
\omega_\gamma[A]=2\langle\alpha,K_0\beta\rangle+R_\omega[A],
\qquad |R_\omega[A]|\le C Q_\gamma[A]^2
\label{eq:endpoint-expansion}
\end{equation}
となる $C>0$ が存在する。
\end{lemma}
\begin{proof}
式\eqref{eq:jacobi}の各 Jacobi 列を $X''=AX$ と書くと、
\[
X(t)=X(0)+tX'(0)+(\mathcal V_A X)(t),\qquad
(\mathcal V_A X)(t):=\int_0^t(t-s)A(s)X(s)\,ds .
\]
区間長が 1 なので
\[
\|\mathcal V_A X\|_\infty
\le \|A\|_{L^1,F}\|X\|_\infty
\le \|A\|_{L^2,F}\|X\|_\infty .
\]
従って $\|A\|_{L^2,F}$ が十分小さければ Neumann--Volterra 級数は絶対収束し、終端の $X(1),X'(1)$、したがって $\Phi(1)$ と $\omega_\gamma[A]$ は $A$ の原点近傍で解析的である。ここで
\[
\mathcal V:=\operatorname{span}\{\sigma_1,\sigma_3\},\qquad
\mathcal E:=\operatorname{span}\{I_2,J_2\}
\]
と置くと、$\mathcal V\mathcal V\subset\mathcal E$ かつ $\mathcal E\mathcal V\subset\mathcal V$ である。従って奇数個の $A(t_i)\in\mathcal V$ の積は再び対称な $\mathcal V$ に入り、反対称チャネルへ寄与しない。一次項も対称なので消え、二次の時間順序項では二つの STF 基底の非可換部分のみが残る。$[\sigma_3,\sigma_1]=-2J_2$ を用いて反対称化すると核 $\frac12(t-s)|t-s|$ が現れ、主項 $2\langle\alpha,K_0\beta\rangle$ を得る。次に寄与しうる次数は四次なので、ある固定近傍で
\[
|R_\omega[A]|\le C_1\|A\|_{L^1,F}^4
\le C_1\|A\|_{L^2,F}^4
=C Q_\gamma^2
\]
と一様に抑えられる。
\end{proof}

\section{鋭い最小作用定理}
\begin{equation}
\kappa_\Omega(\omega):=\inf\{Q_\gamma[A]:\omega_\gamma[A]=\omega\}
\label{eq:value}
\end{equation}
とする。

\begin{theorem}[無拘束の非線形 Jacobi チャネルに対する鋭い法則]
\label{thm:jacobi-sharp}
\begin{equation}
\boxed{\kappa_\Omega(\omega)=\frac{|\omega|}{\|K_0\|}+O(\omega^2),\qquad \omega\to0.}
\label{eq:sharp-jacobi}
\end{equation}
係数は厳密終端問題の意味で鋭い。
\end{theorem}
\begin{proof}
まず Cauchy--Schwarz 不等式と算術・幾何平均不等式から
\begin{equation}
2|\langle\alpha,K_0\beta\rangle|
\le s_0(\|\alpha\|_2^2+\|\beta\|_2^2)=s_0 Q_\gamma.
\label{eq:quad-lower}
\end{equation}
コンパクト性により $-K_0^2$ の最大固有空間はノルムを達成する。実単位ベクトル $\beta_*$ を
\[
-K_0^2\beta_*=s_0^2\beta_*,\qquad
\alpha_*:=K_0\beta_*/s_0
\]
と取ると、$\|\alpha_*\|=\|\beta_*\|=1$ かつ $\langle\alpha_*,K_0\beta_*\rangle=s_0$ である。$A_*=\alpha_*\sigma_3+\beta_*\sigma_1$ とすれば $Q_\gamma[A_*]=2$ かつ二次メモリーは $2s_0$ となり、式\eqref{eq:quad-lower}を飽和する。

厳密回復構成のため $A_\rho=\rho A_*$ とする。前補題より
\begin{equation}
\omega_\gamma[A_\rho]=2s_0 \rho^2+O(\rho^4),\qquad Q_\gamma[A_\rho]=2\rho^2.
\label{eq:recovery-ray}
\end{equation}
解析性と式\eqref{eq:recovery-ray}から
\[
\frac{d}{d\rho}\omega_\gamma[A_\rho]
=4s_0\rho+O(\rho^3)>0
\]
が十分小さい $\rho>0$ で成り立つ。従って $\rho=0$ で通常の逆関数定理を使うのではなく、正の片側での厳密な単調性により、各十分小さい正の $\omega$ に対して一意な $\rho(\omega)>0$ が存在する。従って
\[
Q_\gamma[A_{\rho(\omega)}]=\frac{\omega}{s_0}+O(\omega^2).
\]
負の目標に対しては $A_*^-:=\alpha_*\sigma_3-\beta_*\sigma_1$ と置けば
\[
\omega_\gamma[\rho A_*^-]=-2s_0 \rho^2+O(\rho^4),\qquad
\frac{d}{d\rho}\omega_\gamma[\rho A_*^-]=-4s_0\rho+O(\rho^3)<0,
\]
なので同じ片側単調性で厳密な負の分枝を得る。これが上極限を与える。

一方、上の回復構成からあらかじめ $\kappa_\Omega(\omega)=O(|\omega|)$ が分かるので、最小化列は $Q_\gamma=O(|\omega|)$ を満たすように選べる。その列に式\eqref{eq:endpoint-expansion},\eqref{eq:quad-lower}を適用すると
\[
|\omega|\le s_0 Q_\gamma+C Q_\gamma^2.
\]
$Q_\gamma=O(|\omega|)$ を用いると下界は $Q_\gamma\ge |\omega|/s_0-O(\omega^2)$ となる。回復構成の上界と合わせて式\eqref{eq:sharp-jacobi}が従う。
\end{proof}

Gauss--Legendre Nystr\"om 法による検算では
\begin{equation}
\|K_0\|=0.09103843824\ldots,\qquad \|K_0\|^{-1}=10.98437121\ldots
\label{eq:numeric-k0}
\end{equation}
である。この十進値は数値検証であり、定理の定義そのものではない。

\section{ジャイロスコピック・メモリーとの対応}
\subsection{弱い TT 波の偏光面積}
第2節で固定した局所放射域パッチと適合スクリーン上の同じ先頭放射データを、基準となる大きな Bondi 半径 $r$ での弱い TT ひずみとして
\begin{equation}
h^{TT}(u,r;\eps)=\eps\begin{pmatrix}p(u)&q(u)\\q(u)&-p(u)\end{pmatrix}+O(\eps^2)
\label{eq:tt}
\end{equation}
とする。放射域では $p,q=O(r^{-1})$ である。以下では固定した $r$ で議論し、$r$ 引数を省略する。入射前後で放射が定常なら $\dot p(u_\pm)=\dot q(u_\pm)=0$ である。本節では
\[
X(\eps)=X^{(0)}+\eps X^{(1)}+\eps^2X^{(2)}+O(\eps^3)
\]
と書いたときの \emph{係数}を $X^{(2)}$ と表す。従って $\Delta\Psi_{\rm H}^{(2)}$ は full finite-amplitude memory でも $\eps^2$ を含む寄与そのものでもなく、二次摂動係数である。

計量符号 $(-+++)$、式\eqref{eq:tidal}の $T_{AB}=-R(e_A,k,e_B,k)$、および本稿の outgoing/retarded-time 規約では、局所放射域パッチで $du/d\lambda=1$ と規格化し、スクリーン二脚を Bondi/TT 偏光二脚へ適合させた放射横断ヌル線に対して、先頭の type-N 放射曲率は
\[
R(e_A,k,e_B,k)=-\frac12\ddot h^{TT}_{AB}
+\text{(subleading in $1/r$ and nonlinear amplitude)}
\]
となる。従って先頭次数の潮汐プロファイルは $T=\frac12\ddot h^{TT}$ であり、Jacobi 側と Bondi/TT 側へ同じ $(p,q)$ を入力できる。ここで必要なのはこの固定した局所同定であって、任意の有限 $r$ のヌル方向について同じ $T$ が得られるという主張ではない。TT/Bondi 対応の全体符号は規約依存だが、$(p,q)$ を同時に反転しても以下の二次面積は変わらない。本稿では後述の式\eqref{eq:shear-strain}の符号を固定して全式を一貫させる。

Jacobi チャネルの二次係数は、偏光平面 $(p,q)$ の弦で閉じた符号付き面積
\begin{equation}
\mathcal A_{\rm pol}[p,q]
=\frac12\int_{u_-}^{u_+}(p\dot q-q\dot p)\,du
+\frac12[p_fq_i-p_iq_f]
\label{eq:area}
\end{equation}
に等しく、
\begin{equation}
\omega_\gamma(\eps)=\eps^2\omega_\gamma^{(2)}+O(\eps^3),
\qquad
\omega_\gamma^{(2)}=\mathcal A_{\rm pol}[p,q].
\label{eq:omega-area}
\end{equation}
である。初期シアーなしの基準では $p_i=q_i=0$ なので弦の項は消えるが、最終シアーは非零でよい。

\subsection{Bondi ハード汎関数との係数 2}
局所スクリーン上で
\[
C=\begin{pmatrix}X&Y\\Y&-X\end{pmatrix},\qquad
\widetilde C=\begin{pmatrix}-Y&X\\X&Y\end{pmatrix}
\]
となる。ここでは付録の $\epsilon_{12}=+1$ と、Seraj--Oblak/Faye--Seraj と同じ添字順
$\widetilde C_{AB}=\epsilon_{CA}C^C{}_B$ を用いている。$\epsilon_{AC}$ と添字順を逆にすると全体符号が反転するが、それは本稿の規約ではない。従って
\begin{equation}
N_{AB}\widetilde C^{AB}=2(X\dot Y-Y\dot X).
\label{eq:stf-algebra}
\end{equation}
Faye--Seraj の偏光規約に合わせ、放射域で
\begin{equation}
C_{AB}=-r h^{TT}_{AB}+O(r^0)
\label{eq:shear-strain}
\end{equation}
とする。式\eqref{eq:hard-memory},\eqref{eq:tt},\eqref{eq:stf-algebra}より
\begin{equation}
\Delta\Psi_{\rm H}(\eps)
=\eps^2\Delta\Psi_{\rm H}^{(2)}+O(\eps^3),
\qquad
\Delta\Psi_{\rm H}^{(2)}
=-\frac14\int(p\dot q-q\dot p)\,du.
\label{eq:gyro-area}
\end{equation}
一方、初期シアーなし条件下の式\eqref{eq:omega-area}は
\begin{equation}
\omega_\gamma^{(2)}=\frac12\int(p\dot q-q\dot p)\,du.
\label{eq:jacobi-area-eps}
\end{equation}
従って次を得る。

\begin{theorem}[同一放射データ上の二次汎関数の係数対応]
\label{thm:bridge}
標準的な初期シアーなし Bondi フレーム、弱振幅、放射域の先頭次数、上記の向き付けたスクリーン規約の下で
\begin{equation}
\boxed{\omega_\gamma^{(2)}=-2\Delta\Psi_{\rm H}^{(2)}},\qquad
\boxed{|\omega_\gamma^{(2)}|=2|\Delta\Psi_{\rm H}^{(2)}|}.
\label{eq:bridge}
\end{equation}
\end{theorem}

ここで $\omega_\gamma^{(2)}$ はヌル測地線上の Jacobi 散乱量、$\Delta\Psi_{\rm H}^{(2)}$ は時間的ジャイロ世界線の向き角であり、両者を同一観測量とは主張しない。等式は、第2節で固定した局所放射域パッチ・放射横断ヌル方向・適合スクリーンのもとで、同じ局所 TT/Bondi 放射データ $(p,q)$ を二つの既知汎関数へ代入したときの二次係数の対応である。異なるヌル方向やスクリーン同定へ変更した場合には、その新しい補助プローブに対して対応関係を改めて確認する必要がある。スクリーンの向きを反転すると、両擬スカラーの符号が同時に反転する。有限 $r$ と有限振幅を同時に制御する定理は主張せず、安全な係数の主張は
\begin{equation}
\lim_{r\to\infty}r^2\lim_{\eps\to0}\eps^{-2}
\left(\omega_\gamma(\eps,r)+2\Delta\Psi_{\rm H}(\eps,r)\right)=0
\label{eq:iterated-limit}
\end{equation}
という反復極限として理解する。

\section{入出力定常条件を課した物理的最小作用定理}
入出力の定常性が最良係数を決める。規格化区間で一次潮汐係数を
\[
A^{(1)}(t)=\alpha(t)\sigma_3+\beta(t)\sigma_1,\qquad
\alpha=\frac12p'',\quad \beta=\frac12q''
\]
と書く。このとき $A(t;\eps)=\eps A^{(1)}(t)+O(\eps^2)$ である。固定 $r$ では $p,q=O(r^{-1})$ なので $\alpha,\beta=O(r^{-1})$ であり、$Q_\gamma^{(2)}$ と $\Delta\Psi_{\rm H}^{(2)}$ はともに $O(r^{-2})$ となる。このため以下の比と最良係数には追加の $r$ 因子が残らない。Bel--Robinson 曲率作用は
\begin{equation}
Q_\gamma(\eps)=\eps^2Q_\gamma^{(2)}+O(\eps^3),\qquad
Q_\gamma^{(2)}=\|\alpha\|_2^2+\|\beta\|_2^2.
\label{eq:Q-leading}
\end{equation}
また、$p'(0)=p'(1)=q'(0)=q'(1)=0$ より
\begin{equation}
\int_0^1\alpha(t)\,dt=\int_0^1\beta(t)\,dt=0.
\label{eq:mean-zero}
\end{equation}
したがって
\begin{equation}
\mathcal U_0:=\left\{f\in L^2(0,1):\int_0^1 f(t)\,dt=0\right\},\qquad
P_0f=f-\int_0^1f(s)\,ds
\label{eq:U0}
\end{equation}
を導入し、
\begin{equation}
K_{\rm stat}:=P_0K_0P_0,
\qquad s_{\rm stat}:=\|K_{\rm stat}\|_{\mathcal U_0\to\mathcal U_0}
\label{eq:Kstat}
\end{equation}
とする。$K_{\rm stat}$ は $\mathcal U_0$ 上のコンパクト反自己共役作用素である。

定常条件付きの厳密 Jacobi 最小作用関数を
\begin{equation}
\kappa_{\rm stat}(\omega):=
\inf\{Q_\gamma[A]:\omega_\gamma[A]=\omega,\ \alpha,\beta\in\mathcal U_0\}
\label{eq:kappa-stat}
\end{equation}
と定める。

\begin{theorem}[定常 Jacobi チャネルの鋭い法則]
\label{thm:stationary-sharp}
\begin{equation}
\boxed{\kappa_{\rm stat}(\omega)=\frac{|\omega|}{s_{\rm stat}}+O(\omega^2),\qquad \omega\to0.}
\label{eq:stationary-sharp}
\end{equation}
係数 $s_{\rm stat}^{-1}$ は厳密終端問題の意味で鋭い。
\end{theorem}
\begin{proof}
$\alpha,\beta\in\mathcal U_0$ なら
\[
\langle\alpha,K_0\beta\rangle
=\langle\alpha,P_0K_0P_0\beta\rangle
=\langle\alpha,K_{\rm stat}\beta\rangle.
\]
従って
\[
2|\langle\alpha,K_0\beta\rangle|
\le s_{\rm stat}(\|\alpha\|^2+\|\beta\|^2)
=s_{\rm stat}Q_\gamma.
\]
前節の剰余評価は閉部分空間への制限でも変わらないため、定理\ref{thm:jacobi-sharp}の下界の議論で $s_0$ を $s_{\rm stat}$ に置き換えればよい。

鋭さについては、実 Hilbert 空間 $\mathcal U_0$ 上で $-K_{\rm stat}^2$ の最大固有値 $s_{\rm stat}^2$ に属する実単位固有ベクトル $\beta_*$ を取り、
\[
\alpha_*:=K_{\rm stat}\beta_*/s_{\rm stat}
\]
と置く。すると
\[
\|\alpha_*\|=\|\beta_*\|=1,\qquad
\langle\alpha_*,K_{\rm stat}\beta_*\rangle=s_{\rm stat}
\]
であり、最大特異値対を実関数として構成できる。$A_*=\alpha_*\sigma_3+\beta_*\sigma_1$ とし $A_\rho=\rho A_*$ とすると、平均ゼロ条件は全ての $\rho$ で保存され、
\[
Q_\gamma[A_\rho]=2\rho^2,\qquad
\omega_\gamma[A_\rho]=2s_{\rm stat}\rho^2+O(\rho^4).
\]
定理\ref{thm:jacobi-sharp}と同じ局所反転により正負両方の目標値を厳密に回復できる。よって上極限と下極限が一致し、式\eqref{eq:stationary-sharp}を得る。
\end{proof}

\begin{lemma}[定常 TT ひずみによる実現]
\label{lem:physical-realization}
任意の $\alpha,\beta\in\mathcal U_0$ に対し
\begin{equation}
p(t)=2\int_0^t(t-s)\alpha(s)\,ds,\qquad
q(t)=2\int_0^t(t-s)\beta(s)\,ds
\label{eq:pq-from-control}
\end{equation}
と置けば
\[
p(0)=q(0)=p'(0)=q'(0)=p'(1)=q'(1)=0,
\qquad p''=2\alpha,\ q''=2\beta.
\]
従って、初期シアーなし、かつ入出力定常な一次 TT 波形を与える。
\end{lemma}
\begin{proof}
初期値は定義から明らかであり、$p'(t)=2\int_0^t\alpha(s)ds$ だから $p'(1)=2\int_0^1\alpha=0$ である。$q$ も同様である。最終値 $p(1),q(1)$ は一般に非零であり、一定の最終シアーを許す。
\end{proof}

\begin{lemma}[滑らかな定常バーストでも同じ鋭い定数]
\label{lem:smooth-density}
\[
\mathcal U_{0,c}^\infty:=\left\{f\in C_c^\infty(0,1):\int_0^1f(t)\,dt=0\right\}
\]
は $\mathcal U_0$ に稠密である。従って、曲率が端点近傍で消える滑らかな定常 TT バーストへ許容クラスを狭めても、先頭次数の鋭い係数 $s_{\rm stat}$ は変わらない。
\end{lemma}
\begin{proof}
$f\in\mathcal U_0$ に対し $g_n\in C_c^\infty(0,1)$ を $L^2$ で $f$ に近づける。$\int g_n\to0$ である。$\int\chi=1$ を満たす固定 $\chi\in C_c^\infty(0,1)$ を取り、
\[
f_n:=g_n-\left(\int_0^1g_n\right)\chi
\]
と置けば $f_n\in\mathcal U_{0,c}^\infty$ かつ $f_n\to f$ in $L^2$ である。$K_{\rm stat}$ は有界なので双線形形式 $\langle\alpha,K_{\rm stat}\beta\rangle$ は $L^2\times L^2$ で連続である。従って最大特異値対をこの滑らかな平均ゼロ関数で任意精度に近似できる。さらに各近似対は端点近傍で曲率が消えるため、式\eqref{eq:pq-from-control}で得た $p,q$ は区間外へ定数として滑らかに接続できる。よって smooth stationary burst class の infimum も同じ鋭い定数を持つ。
\end{proof}

\section{主定理：Bel--Robinson 曲率曝露の鋭い下界}
先頭次数のハード・メモリー目標 $\psi=\Delta\Psi_{\rm H}^{(2)}$ に対し、\emph{leading-order 放射データ問題}の値関数を
\begin{equation}
\kappa_{\rm H}^{\rm lead}(\psi):=
\inf\left\{Q_\gamma^{(2)}[A^{(1)}]:\Delta\Psi_{\rm H}^{(2)}[A^{(1)}]=\psi,\ \alpha,\beta\in\mathcal U_0\right\}
\label{eq:kappa-hard-lead}
\end{equation}
と定める。ここで最小化は上記の弱振幅 TT/Bondi 接線データのクラスで行い、full finite-amplitude Einstein 解全体の値関数は定義しない。

\begin{theorem}[固定放射横断ヌル窓に対するハード・メモリーの鋭い最小曲率作用]
\label{thm:main-physical}
第2節の Bel--Robinson 規約、$L/r\to0$ の局所放射域パッチ内で固定した放射横断ヌル窓 $\gamma$ と適合スクリーン、標準的な初期シアーなし Bondi フレーム、弱振幅かつ放射域の先頭次数、入出力定常条件の下で（最終シアーは零に固定しない）
\begin{equation}
\boxed{\kappa_{\rm H}^{\rm lead}(\psi)=\frac{2}{\|K_{\rm stat}\|}|\psi|.}
\label{eq:main-physical-value}
\end{equation}
この等式は摂動係数としての先頭次数放射データ問題で厳密であり、$2/\|K_{\rm stat}\|$ は最良係数である。
\end{theorem}
\begin{proof}
式\eqref{eq:endpoint-expansion}の二次主項と定理\ref{thm:bridge}から、初期シアーなし・定常クラスでは
\[
\Delta\Psi_{\rm H}^{(2)}=-\langle\alpha,K_{\rm stat}\beta\rangle .
\]
従って
\[
2|\psi|=2|\langle\alpha,K_{\rm stat}\beta\rangle|
\le s_{\rm stat}(\|\alpha\|_2^2+\|\beta\|_2^2)=s_{\rm stat}Q_\gamma^{(2)},
\]
すなわち $Q_\gamma^{(2)}\ge2|\psi|/s_{\rm stat}$ である。

逆に、定理\ref{thm:stationary-sharp}の証明で構成した実最大特異値対 $\alpha_*,\beta_*$ を用いる。目標の符号に応じて $\beta_*$ の符号を選び、
\[
\alpha=r\alpha_*,\qquad \beta=\pm r\beta_*,\qquad r^2=\frac{|\psi|}{s_{\rm stat}}
\]
と置けば、$\Delta\Psi_{\rm H}^{(2)}=\psi$ かつ
\[
Q_\gamma^{(2)}=2r^2=\frac{2}{s_{\rm stat}}|\psi|
\]
を厳密に満たす。補題\ref{lem:physical-realization}で初期シアーなし・入出力定常な TT データへ持ち上がり、補題\ref{lem:smooth-density}により端点近傍で曲率が消える滑らかなバーストに限定しても同じ infimum を得る。
\end{proof}

従って固定ヌル窓上の下界は先頭次数の限定クラスで
\begin{equation}
Q_\gamma^{(2)}\ge\frac{2}{\|K_{\rm stat}\|}|\Delta\Psi_{\rm H}^{(2)}|
\label{eq:main-bound}
\end{equation}
と書ける。数値検証では
\begin{equation}
\|K_{\rm stat}\|=0.0281140680170\ldots,
\qquad
\boxed{\frac{2}{\|K_{\rm stat}\|}=71.13876223\ldots.}
\label{eq:main-bound-num}
\end{equation}
を得る。この十進値は定理の定義ではなく数値検証である。

\begin{remark}[旧係数 $21.96874242\ldots$ との関係]
無拘束 $L^2$ 問題の係数 $2/\|K_0\|=21.96874242\ldots$ は、定常部分クラスに対しても有効な弱い下界である。しかし、定常クラスの最良係数ではない。平均ゼロ制約の下では最適化を $K_{\rm stat}=P_0K_0P_0$ 上でやり直す必要があり、その結果が上の係数 $2/\|K_{\rm stat}\|$ である。
\end{remark}

\section{フレーム依存性と適用範囲}
\subsection{初期シアーなしフレームの必要性}
$u$ に依存しないシアーのずれ
\[
C_{AB}\mapsto C_{AB}+S_{AB}
\]
ではニュースは変わらないが、ハード汎関数は
\[
\bar\Phi_{\rm H}\mapsto\bar\Phi_{\rm H}
-\frac18\Delta C_{AB}\widetilde S^{AB}
\]
だけ変化する。一方、式\eqref{eq:area}の弦で閉じた面積は $(p,q)\mapsto(p+p_0,q+q_0)$ に不変である。したがって、Jacobi 面積とハード項単独の同一視は一般の一定シアー基準では成立しない。初期シアーなしフレームは単なる記法上の便宜ではなく、ハード成分だけの対応関係を成立させる物理的規約である。

\subsection{固定ヌル窓が補助構造であること}
$Q_\gamma$ は波形だけから決まる固有量ではなく、固定したヌル窓 $\gamma$ を含む補助的なコストである。例えば曲率の台と放射履歴を変えずに、両端へ曲率ゼロ・ニュースゼロの定常区間だけを付加し、同じ $du/d\lambda=1$ 規格化のまま窓長を $L$ から $L'>L$ へ延長する。このとき $\int_\gamma B(k^4)d\lambda$ とハード・メモリーは変わらない一方、
\[
Q_{\gamma'}=\left(\frac{L'}{L}\right)^3Q_\gamma
\]
となる。従って同じ物理波形について窓を延長すれば $Q_\gamma$ を任意に大きくできる。本稿の最小作用則は、固定窓ごとの無次元化された変分問題に対するものであり、異なる窓の $Q_\gamma$ を同一の detector exposure として比較する主張ではない。なお時間軸と波形全体を同時にアフィン再パラメータ化する場合は、第2節で示した通り $Q_\gamma$ は不変であり、これは空の定常尾を追加して物理的な窓そのものを変える操作とは異なる。

\subsection{定常条件が最良係数を変える理由}
無拘束最適化では $K_0$ の最大特異値部分空間が回復族を与えるが、定常な放射では $\int\alpha=\int\beta=0$ が必要である。この条件は各偏光に余次元 1 の閉部分空間を課すため、稠密近似によって無視できない。実際、全空間の最大特異値部分空間に $\int\beta=0$ と $\int K_0\beta=0$ を同時に課す線形系は数値的に非退化であり、旧最適化波形は定常クラスに入らない。正しい接線方向作用素は $K_{\rm stat}=P_0K_0P_0$ であり、$\|K_{\rm stat}\|<\|K_0\|$ なので、必要な鋭い曲率曝露係数は増加する。

本稿は最終シアーを零には固定しない。もしさらに $p(1)=q(1)=0$ を課すなら、平均ゼロ条件に加えて
\[
\int_0^1(1-t)\alpha(t)\,dt=0,\qquad
\int_0^1(1-t)\beta(t)\,dt=0
\]
という一次モーメント制約が加わる。その場合の許容部分空間と圧縮作用素は $K_{\rm stat}$ とは異なり、最良係数も一般には変わりうるため、本稿の $71.13876223\ldots$ をその追加制約付き問題へ移してはならない。

\subsection{主張しないもの}
本稿は次を主張しない。
\begin{itemize}
\item ジャイロスコピック・メモリー全体への鋭い下界；
\item スピン・メモリーへの鋭い下界；
\item 一般の BMS フレームで不変なハード成分のみの定理；
\item 有限振幅での大域的最小作用則；
\item 本稿の回復 TT/Bondi データが任意に完全な漸近平坦真空 Einstein 時空へ大域実現できるという主張；
\item 有限 $r$ での全次数の厳密恒等式；
\item Bondi エネルギー、Isaacson エネルギー、放射エネルギー、検出器エネルギーの下界；
\item ジャイロスコピック・メモリー、Bel--Robinson テンソル、Jacobi メモリーテンソル自体の発見。
\end{itemize}

従って主定理\ref{thm:main-physical}の「物理クラス」は、完全 Einstein 解の大域解空間ではなく、標準的な初期シアーなし規約のもとで整合する弱振幅・放射域先頭次数の TT/Bondi 放射データクラスを意味する。比較している量も full finite-amplitude の $Q_\gamma$ と $\Delta\Psi_{\rm H}$ ではなく、その二次摂動係数 $Q_\gamma^{(2)}$ と $\Delta\Psi_{\rm H}^{(2)}$ である。完全な非線形時空への実現可能性と有限振幅での値関数は別問題として残す。

\section{数値検証と再現性}
解析証明と数値検証を区別する。式\eqref{eq:sharp-jacobi},\eqref{eq:stationary-sharp},\eqref{eq:bridge},\eqref{eq:main-physical-value}の論理自体は解析的であり、以下の数値は規格化、符号、終端量の抽出、制限スペクトル値、および反例志向の検査例を確認するためのものである。

\subsection{\texorpdfstring{無拘束 $K_0$ と厳密 Jacobi 回復}{無拘束 K0 と厳密 Jacobi 回復}}
Gauss--Legendre Nystr\"om 離散化は
\begin{center}
\begin{tabular}{rr}
\toprule
節点数 & $\|K_0\|$\\
\midrule
50&0.0910384416041\\
100&0.0910384384527\\
200&0.0910384382516\\
300&0.0910384382407\\
500&0.0910384382384\\
\bottomrule
\end{tabular}
\end{center}
を与える。同じ離散最大特異値対を厳密 Jacobi ODE に入れると
\begin{center}
\begin{tabular}{rr}
\toprule
$\rho$ & $Q_\gamma/|\omega_\gamma|$\\
\midrule
0.03&10.98435328\\
0.01&10.98436922\\
0.003&10.98437103\\
0.001&10.98437118\\
\bottomrule
\end{tabular}
\end{center}
となり、$\|K_0\|^{-1}=10.98437121\ldots$ に収束する。ここで $\rho$ は回復半直線 $A_\rho=\rho A_*$ の振幅であり、物理摂動パラメータ $\eps$ とは別である。反対称 Jacobi チャネルでは奇数次数が消えるため $\omega_\gamma[A_\rho]=2s_0\rho^2+O(\rho^4)$ となり、比の誤差が $O(\rho^2)$ で収束することとも整合する。

\subsection{定常条件付きスペクトルと厳密回復}
平均ゼロ射影を離散化して $P_0MP_0$ の最大特異値を計算すると、
\begin{center}
\begin{tabular}{rrr}
\toprule
節点数&$\|K_{\rm stat}\|$&$2/\|K_{\rm stat}\|$\\
\midrule
50&0.0281140528760&71.13880054294\\
100&0.0281140670547&71.13876466581\\
200&0.0281140679578&71.13876238062\\
300&0.0281140680065&71.13876225721\\
500&0.0281140680170&71.13876223067\\
\bottomrule
\end{tabular}
\end{center}
を得る。さらに連続な平均ゼロ最大特異値対を厳密 Jacobi ODE へ入れると
\begin{center}
\begin{tabular}{rr}
\toprule
$\rho$ & $Q_\gamma/|\omega_\gamma|$\\
\midrule
0.03&35.56942171300\\
0.01&35.56938562391\\
0.003&35.56938151880\\
0.001&35.56938115793\\
\bottomrule
\end{tabular}
\end{center}
となり、$\|K_{\rm stat}\|^{-1}=35.5693811\ldots$ へ収束する。ここでも比の誤差は上記の偶数次数構造により $O(\rho^2)$ である。再構成したひずみは数値精度で $p'(1)=q'(1)=0$ を満たす。連続最大特異値対を $\|\alpha_*\|=\|\beta_*\|=1$ と規格化すると、数値 SVD の一つの向き付けでは $p_f/\rho\simeq-0.47863$, $q_f/\rho\simeq0.30730$ である。最大特異ベクトル対には全体符号の任意性があるため、両者を同時に反転した $(0.47863,-0.30730)$ も同じ最適化器を表す。偏光基底の回転とこの全体符号に依存しない最終シアーの大きさは $\sqrt{p_f^2+q_f^2}/\rho\simeq0.56878$ である。本定理は最終シアーを零に固定せず、一定の displacement-memory 成分を許している。これらは制限定理そのものの証明ではなく、解析的回復構成の実装検査である。

\subsection{対応関係の波形検査}
$W=\int(p\dot q-q\dot p)du$ とすると、初期ゼロ基準の場合 $\omega^{(2)}=W/2$, $\Delta\Psi_{\rm H}^{(2)}=-W/4$ である。代表的な検査例は
\begin{center}
\begin{tabular}{lrrrr}
\toprule
波形族&$W$&$\omega^{(2)}$&$\Delta\Psi_{\rm H}^{(2)}$&比\\
\midrule
滑らかな円偏光&2.35619447&1.17809724&-0.58904862&-2.000000\\
逆向きヘリシティ&-2.35619447&-1.17809724&0.58904862&-2.000000\\
楕円偏光&1.22522112&0.61261056&-0.30630528&-2.000000\\
固定軸&$\simeq0$&$\simeq0$&$\simeq0$&ともにゼロ\\
終端メモリー&-0.62590250&-0.31295125&0.15647563&-2.000000\\
ヘリシティ相殺&$\simeq0$&$\simeq0$&$\simeq0$&ともにゼロ\\
\bottomrule
\end{tabular}
\end{center}
である。基準値に対する反例志向の検査では、終端メモリー波形に $(p_0,q_0)=(0.35,-0.22)$ を加える前後で、弦で閉じた面積は $-0.3129512514\to-0.3129512515$ と不変だが、ハード・ジャイロ汎関数は $0.1564756257\to0.1124756258$ と変化する。これは一般の基準値に対するハード成分のみの恒等式を反証し、初期シアーなしフレーム制限を支持する。

\subsection{旧最適化波形の不適合を検出する監査}
全空間の $K_0$ 最大特異値部分空間に正規直交基底を取り、$\int\beta=0$ と $\int K_0\beta=0$ の二つの制約行を $\|K_0\|$ で規格化すると、その $2\times2$ 行列式の絶対値は基底の直交変換では不変で、$n=500$ で約 $0.90925$ となる。この量は同時に、定数関数の最大特異部分空間への直交射影のノルム二乗に等しい。従って定数モードの約 $90.9\%$ が旧最大特異平面に含まれ、二つの定常制約がその平面上で独立であることを示す。これは旧無拘束最大特異値部分空間が定常条件を自動的には満たさないことを示す。この監査は「なぜ旧係数が鋭くなかったか」を診断するためのものであり、新係数の鋭さ自体は定理\ref{thm:stationary-sharp}の解析証明から従う。

\section{先行研究との境界と新規性の表現}
本稿の構成要素を、既知の内容、既知結果から直接従う内容、本稿で導く結果に分ける。
\begin{center}
\begin{tabular}{p{0.48\linewidth}p{0.42\linewidth}}
\toprule
要素&本稿での分類\\
\midrule
ジャイロスコピック重力メモリー&既知\cite{SerajOblak2023}\\
ソフト／スピン成分とハード／ヘリシティ成分の分解&既知\cite{SerajOblak2023,FayeSeraj2025}\\
$N_{AB}\widetilde C^{AB}$ 型ハード汎関数&既知\cite{SerajOblak2023,FayeSeraj2025}\\
ヘリシティ／局所双対性生成子としての解釈&既知\cite{SerajOblak2023}\\
Jacobi に基づく持続的観測量／非線形平面波&既知\cite{Flanagan2019,Flanagan2020}\\
平面波メモリーテンソル&既知\cite{Harte2025}\\
Bondi--Sachs における曲線偏差観測量の階層&既知\cite{GrantNichols2022}\\
ホロノミーから導くジャイロスコピック・メモリー&既知\cite{SerajNeogi2023}\\
Bel--Robinson 曲率代数&既知\cite{FerrandoSaez2012}\\
面積制約付き曲線最適化・曲げエネルギー&隣接する既知の変分問題\cite{Hahmann2005,Ferone2016}\\
偏光面積とハード・ジャイロ汎関数の係数対応&既知汎関数からの規約固定した導出\\
無拘束 Jacobi 最小曝露則 $1/\|K_0\|$&本稿で導出\\
定常 Jacobi 最小作用則 $1/\|P_0K_0P_0\|$&本稿で導出\\
固定ヌル窓付きハード・メモリー問題の最良係数 $2/\|P_0K_0P_0\|$&本稿で導出\\
\bottomrule
\end{tabular}
\end{center}

Flanagan--Grant--Harte--Nichols は、一般の持続的観測量と非線形平面波における Jacobi 構造を確立している\cite{Flanagan2019,Flanagan2020}。Grant--Nichols は、漸近平坦時空における曲線偏差観測量を Bondi--Sachs 枠組みで展開した\cite{GrantNichols2022}。Harte--Mieling--Oancea--Steininger は、平面波散乱を少数のメモリーテンソルで整理している\cite{Harte2025}。Seraj--Neogi は、漸近ホロノミーの Lorentz 部分とジャイロスコピック・メモリーの関係を示した\cite{SerajNeogi2023}。これらを本稿独自の構造として数えない。

数学側にも近縁問題がある。Hahmann--Sauvage--Bonneau は、閉平面曲線の面積を保つ multiresolution deformation と、その最適化で扱う曲げエネルギーを研究している\cite{Hahmann2005}。Ferone--Kawohl--Nitsch は、固定面積を持つ単純閉曲線について幾何学的な elastica 最小化を研究している\cite{Ferone2016}。これらは本稿の符号付き面積汎関数と二階微分コストに近いが、閉曲線／弧長幾何、許容境界条件、固定ヌル窓の Jacobi/Bondi 対応、平均ゼロ圧縮作用素という本稿の同条件問題とは異なる。従って、本稿は「面積制約付き曲げエネルギー最適化」一般を新規とは主張しない。

本稿は歴史的優先権を主張しない。また、双線形形式に対する Cauchy--Schwarz 不等式、コンパクト作用素ノルムの飽和、面積制約付き曲線最適化という一般的機構それ自体は既知であり、そこを新規性とは数えない。最終定理の同時条件について、入出力定常、初期シアーなし、固定した放射横断ヌル窓、Bel--Robinson 曲率作用、ハード・ジャイロ／ヘリシティ目標、最良係数を同時に扱う先行定理は文献検索で特定できなかった。本稿が新規性として主張する範囲は、物理的な入出力定常条件から平均ゼロ制約と $K_{\rm stat}=P_0K_0P_0$ を特定し、固定した補助ヌルプローブ上の曲率作用とハード・ジャイロ汎関数を同じ放射データ上の制約付き逆最適化問題として結び、exact Jacobi 回復まで含めて最良係数を確定する組合せに限定する。この文献検索結果は歴史的な世界初を保証するものではない。

\section{考察}
本稿の主結果は、確立されたハード・ジャイロ汎関数と、同じ放射データを横切る固定ヌル窓上の曲率作用の先頭摂動係数との間に最良下界を与える。ハード・ジャイロスコピック・メモリーの二次係数を指定すると、それより小さな $Q_\gamma^{(2)}$ で同じ先頭次数のハード／ヘリシティ・チャネルを生成することはできず、さらに定常条件を満たす滑らかな許容回復族がその境界へ任意に近づく。これは既知のヘリシティ汎関数の単なる記号置換ではなく、波形空間上の制約付き逆最適化を伴う。ただし、最適化の最終段階そのものは標準的な作用素ノルム問題であり、結果の学術的価値はその物理的な問題設定・制約作用素・対応関係が同条件で既知かどうかに依存する。

定常性が係数を大きく変える点も重要である。無拘束 Jacobi 最小作用問題の係数はハード・メモリーへ形式的に移せるものの、その最適化波形はバーストの入出力定常性を保証しない。物理条件を正しく課すと、許容接空間は $\mathcal U_0$ に縮み、作用素ノルムは $\|K_0\|$ から $\|P_0K_0P_0\|$ に変わる。従って最良係数は約 $21.97$ ではなく約 $71.14$ となる。この差は平均ゼロ制約による作用素圧縮そのものに由来する。

一方で射程は限定される。第一に $Q_\gamma$ は標準的な重力波エネルギーではない。第二にハード・ジャイロスコピック・メモリーは $r^{-2}$ かつ弱重力場では二次であり、現時点で検出器に直接結びつく観測量としては弱い。第三にハード／ソフト分解はフレームに依存し、本稿は初期シアーなし規約に固定される。第四に本定理は弱振幅・放射域の先頭次数に関する主張であり、有限振幅における大域的な非線形時空の最小作用関数までは扱わない。

従って本稿の結論は三層である。(i) 非線形 Jacobi メモリー・チャネルには、小目標に対する厳密で鋭い最小曝露則がある。(ii) 入出力定常を課すと、同じ構造は平均ゼロ部分空間へ圧縮した作用素 $K_{\rm stat}=P_0K_0P_0$ 上で、再び厳密で鋭い法則を持つ。(iii) 標準的な初期シアーなし Bondi 放射では、その定常 Jacobi チャネルの二次係数はハード・ジャイロスコピック・メモリーの二次係数の二倍であり、$Q_\gamma^{(2)}$ をコストとする leading-order 放射データ問題は $\kappa_{\rm H}^{\rm lead}(\psi)=2|\psi|/\|K_{\rm stat}\|$ と厳密に解ける。

\section{結論}
本稿ではヌル Jacobi 散乱に対する最小作用問題を自己完結に定式化し、まず無拘束クラスで
\[
\kappa_\Omega(\omega)=\frac{|\omega|}{\|K_0\|}+O(\omega^2)
\]
を証明した。次に、物理的な入出力定常性が潮汐制御の平均ゼロ条件を課すことを用い、制限作用素 $K_{\rm stat}=P_0K_0P_0$ に対して
\[
\kappa_{\rm stat}(\omega)=\frac{|\omega|}{\|K_{\rm stat}\|}+O(\omega^2)
\]
という鋭い法則を得た。さらに Bondi シアーの局所 STF 代数と弱い TT 波の偏光面積を通じて、$\omega_\gamma^{(2)}=-2\Delta\Psi_{\rm H}^{(2)}$ を導いた。従って、標準的な初期シアーなしフレーム、入出力定常、弱振幅かつ放射域の先頭次数では
\[
\kappa_{\rm H}^{\rm lead}(\psi)=\frac{2}{\|K_{\rm stat}\|}|\psi|,
\qquad
\frac{2}{\|K_{\rm stat}\|}=71.13876223\ldots
\]
である。ここで $\psi=\Delta\Psi_{\rm H}^{(2)}$ は二次摂動係数であり、最小化されるコストは $Q_\gamma^{(2)}$ である。この係数は leading-order TT/Bondi 放射データ問題で厳密に鋭い。一方、exact nonlinear Jacobi endpoint の値関数は
$\kappa_{\rm stat}(\omega)=|\omega|/\|K_{\rm stat}\|+O(\omega^2)$
という別の漸近定理であり、full finite-amplitude Einstein 時空におけるハード・メモリー最小作用則は本稿の主張外である。

今後の理論的課題として、固定ヌル窓という補助構造を除去できるか、有限振幅の Einstein 解空間へ値関数を拡張できるか、最終シアーにも制約を課した場合に最良係数がどう変わるかが残る。

\appendix
\section{規約一覧}
\begin{center}
\begin{tabular}{p{0.32\linewidth}p{0.58\linewidth}}
\toprule
項目&規約\\
\midrule
計量の符号&$(-+++)$\\
Bondi ニュース&$N_{AB}=\partial_u C_{AB}$\\
初期フレーム&$C_{AB}(-\infty)=0$\\
スクリーンの向き&$\epsilon_{12}=+1$\\
双対シアー&$\widetilde C_{AB}=\epsilon_{CA}C^C{}_B$\\
TT/Bondi 対応&$C_{AB}=-rh^{TT}_{AB}+O(r^0)$（偏光規約を固定）\\
光学的潮汐行列&$T_{AB}=-R(e_A,k,e_B,k)$\\
Bel--Robinson テンソル&式\eqref{eq:BR-def}、明示的な $1/4$ の係数\\
曲率曝露&$Q_\gamma=L^3\int B(k^4)d\lambda$（$du/d\lambda=1$ の代表では $d\lambda=du$）\\
先頭物理コスト&$Q_\gamma(\eps)=\eps^2Q_\gamma^{(2)}+O(\eps^3)$\\
非線形 Jacobi スカラー&$\operatorname{skew}(S_{11}+S_{22})=\omega_\gamma J_2$\\
定常制御空間&$\mathcal U_0=\{f\in L^2(0,1):\int f=0\}$\\
\bottomrule
\end{tabular}
\end{center}

\section{\texorpdfstring{$K_0$ と制限作用素のスペクトル構造}{K0 と制限作用素のスペクトル構造}}
$K_0$ はコンパクト反自己共役作用素であり、$-K_0^2$ は正のコンパクト自己共役作用素である。最大固有値 $s^2$ は $\|K_0\|^2$ に等しい。また分布微分または弱微分の意味で
\[
(K_0f)'''=2f
\]
なので、固有方程式を境界条件を伴う ODE 系に変換できる。本稿では厳密な定数を作用素ノルムで定義し、十進小数は証明に用いない。

定常射影について、$k_0(t,s)=\frac12(t-s)|t-s|$ と
\[
m(t)=\int_0^1k_0(t,s)ds=\frac{t^3-(1-t)^3}{6}
\]
を置くと、$P_0K_0P_0$ は核
\begin{equation}
k_{\rm stat}(t,s)=k_0(t,s)-m(t)+m(s)
\label{eq:kstat-kernel}
\end{equation}
を持つ。これは $k_{\rm stat}(s,t)=-k_{\rm stat}(t,s)$ であり二乗可積分なので、$K_{\rm stat}$ は $\mathcal U_0$ 上のコンパクト反自己共役作用素である。さらに $(K_{\rm stat}f)'''=2f$ だから $K_{\rm stat}\neq0$。従って $-K_{\rm stat}^2$ の最大正固有値が存在し、その平方根が $s_{\rm stat}=\|K_{\rm stat}\|$ である。

Gauss--Legendre 節点 $t_i$ と重み $w_i$ に対する Nystr\"om 行列
\[
M_{ij}=\sqrt{w_i}\,\frac12(t_i-t_j)|t_i-t_j|\,\sqrt{w_j}
\]
を用いる。定数関数に対応する単位ベクトル $v_i=\sqrt{w_i}$ に対して $P=I-vv^T$ と置くと、$PMP$ の最大特異値が $\|K_{\rm stat}\|$ の浮動小数点近似を与える。

\section{厳密回復構成と定常波形による実現の詳細}
定理\ref{thm:jacobi-sharp}の無拘束回復構成は、二次最適化波形を厳密終端写像上の半直線へ持ち上げる。 STF 制御行列の積では奇数次数の積は対称部分空間に戻るため、反対称終端チャネルの三次項は消え、二次項の次は四次項である。$A_\rho=\rho A_*$ に沿って $Q=2\rho^2$ は厳密で、$\omega=2s_0 \rho^2+O(\rho^4)$ である。終端写像は $\rho$ に解析的なので
\[
\frac{d\omega}{d\rho}=4s_0\rho+O(\rho^3)>0
\]
が十分小さい $\rho>0$ で成り立つ。従って $\rho=0$ で微分が消える通常の逆関数定理には依存せず、正の片側単調性から近似目標ではない厳密な $\omega$ を実現する $\rho(\omega)$ が存在する。$\beta_*$ を反転した分枝では導関数が負になり、小さい負の目標も同様に厳密回復される。

定常問題も同じ構造を持つが、最適化波形を $K_{\rm stat}$ の最大特異値対から取る点だけが異なる。実単位固有ベクトル $-K_{\rm stat}^2\beta_*=s_{\rm stat}^2\beta_*$ から $\alpha_*=K_{\rm stat}\beta_*/s_{\rm stat}$ と構成できるので、対は実関数として選べる。$\alpha_*,\beta_*\in\mathcal U_0$ なので、半直線 $\rho A_*$ は全振幅で接線方向の定常制約を保存する。また
\[
p_\rho(t)=2\rho\int_0^t(t-s)\alpha_*(s)ds,\qquad
q_\rho(t)=2\rho\int_0^t(t-s)\beta_*(s)ds
\]
は
\[
p_\rho(0)=q_\rho(0)=p'_\rho(0)=q'_\rho(0)=p'_\rho(1)=q'_\rho(1)=0
\]
を満たす。従って、初期ゼロ基準、入出力定常、任意の一定最終シアーを同時に許す物理的な接線方向の波形履歴である。これが、無拘束クラスでの鋭さから物理クラスでの鋭さへ移す際に不足していた許容性の確認を閉じる。

\section{Bel--Robinson 規格化の確認}
式\eqref{eq:BR-def}の $1/4$ 規約を固定する。真空のヌル・スクリーンでは Weyl 潮汐データが二つの実自由度に縮退する。$m=(e_1+i e_2)/\sqrt2$ と置くと、$T=\alpha\sigma_3+\beta\sigma_1$ に対して $\Psi_0=C(k,m,k,m)$ は本稿の符号規約では全体符号を除き $\alpha+i\beta$ に一致する。従って
\[
|\Psi_0|^2=\alpha^2+\beta^2=\frac12\tr(T^2).
\]
式\eqref{eq:BR-def}の明示的な $1/4$ 規約では四重ヌル縮約が $B(k,k,k,k)=|\Psi_0|^2$ となるため、
\[
B(k,k,k,k)=|\Psi_0|^2=\frac12\tr(T^2).
\]
これより式\eqref{eq:Q},\eqref{eq:control}が従う。異なる Bel--Robinson 規格化を使う文献とは係数が変わるため、比較時にはこの $1/4$ を固定する。

\section{摂動次数と極限の順序}
Bondi/Jacobi の対応関係は放射域の先頭次数に関する主張である。$C_{AB}=\eps C_{AB}^{(1)}+O(\eps^2)$, $h_{AB}^{TT}=-C_{AB}/r+O(r^{-2})$ とする。本稿の $X^{(2)}$ は
\[
X(\eps)=X^{(0)}+\eps X^{(1)}+\eps^2X^{(2)}+O(\eps^3)
\]
における二次摂動\emph{係数}を表す。特に
\[
Q_\gamma(\eps)=\eps^2Q_\gamma^{(2)}+O(\eps^3),\quad
\omega_\gamma(\eps)=\eps^2\omega_\gamma^{(2)}+O(\eps^3),\quad
\Delta\Psi_{\rm H}(\eps)=\eps^2\Delta\Psi_{\rm H}^{(2)}+O(\eps^3).
\]
式\eqref{eq:iterated-limit}の内側の $\eps^{-2}$ 極限はこれらの二次係数を抽出し、その後 $r\to\infty$ を取る。従って、$O(\eps^3)$ と $O(r^{-3})$ の剰余を一つの一様定数で同時制御する有限 $r$・有限振幅定理は主張しない。

\section{再現性一覧と定常性監査}
解析証明と浮動小数点による数値検証を区別する。数値計算は、全空間の $K_0$ 最適化波形が平均ゼロの定常制約を自動的には満たさないこと、$P_0K_0P_0$ の Nystr\"om 収束、および最大特異値対を用いた小振幅での厳密 Jacobi 回復を検算する。最良係数の論理はコンパクト作用素による解析から成立し、浮動小数点値 $71.13876223\ldots$ は数値検証としてのみ用いる。

\begin{thebibliography}{99}
\bibitem{Flanagan2019}
E. E. Flanagan, A. M. Grant, A. I. Harte, and D. A. Nichols,
``Persistent gravitational wave observables: General framework,''
Phys. Rev. D 99, 084044 (2019), arXiv:1901.00021, DOI: 10.1103/PhysRevD.99.084044.

\bibitem{Flanagan2020}
E. E. Flanagan, A. M. Grant, A. I. Harte, and D. A. Nichols,
``Persistent gravitational wave observables: Nonlinear plane wave spacetimes,''
Phys. Rev. D 101, 104033 (2020), arXiv:1912.13449, DOI: 10.1103/PhysRevD.101.104033.

\bibitem{GrantNichols2022}
A. M. Grant and D. A. Nichols,
``Persistent gravitational wave observables: Curve deviation in asymptotically flat spacetimes,''
Phys. Rev. D 105, 024056 (2022), arXiv:2109.03832, DOI: 10.1103/PhysRevD.105.024056; Erratum Phys. Rev. D 107, 109902 (2023).

\bibitem{SerajOblak2023}
A. Seraj and B. Oblak,
``Gyroscopic Gravitational Memory,''
JHEP 11, 057 (2023), arXiv:2112.04535, DOI: 10.1007/JHEP11(2023)057.

\bibitem{SerajNeogi2023}
A. Seraj and T. Neogi,
``Memory effects from holonomies,''
Phys. Rev. D 107, 104034 (2023), arXiv:2206.14110, DOI: 10.1103/PhysRevD.107.104034.

\bibitem{FayeSeraj2025}
G. Faye and A. Seraj,
``Gyroscopic Gravitational Memory from quasi-circular binary systems,''
Class. Quantum Grav. 42, 035005 (2025), arXiv:2409.02624, DOI: 10.1088/1361-6382/ada339.

\bibitem{Harte2025}
A. I. Harte, T. B. Mieling, M. A. Oancea, and E. Steininger,
``Gravitational wave memory and its effects on particles and fields,''
Phys. Rev. D 111, 024034 (2025), arXiv:2407.00174, DOI: 10.1103/PhysRevD.111.024034.

\bibitem{FerrandoSaez2012}
J. J. Ferrando and J. A. S\'aez,
``On the Bel radiative gravitational fields,''
Class. Quantum Grav. 29, 075012 (2012), arXiv:1111.2755, DOI: 10.1088/0264-9381/29/7/075012.

\bibitem{Hahmann2005}
S. Hahmann, B. Sauvage, and G.-P. Bonneau,
``Area preserving deformation of multiresolution curves,''
Comput. Aided Geom. Des. 22, 349--367 (2005), DOI: 10.1016/j.cagd.2005.01.006.

\bibitem{Ferone2016}
V. Ferone, B. Kawohl, and C. Nitsch,
``The elastica problem under area constraint,''
Math. Ann. 365, 987--1015 (2016), arXiv:1411.6100, DOI: 10.1007/s00208-015-1284-y.
\end{thebibliography}

\end{document}
